\chapter*{Abstract}
\addcontentsline{toc}{chapter}{ABSTRACT}

Automated facial analysis systems are widely deployed despite well-documented accuracy disparities across demographic subgroups, and independent, reproducible auditing of these systems remains rare in practice. This report proposes BiasAperture, a diagnostic and evaluative software platform that computes subgroup and intersectional fairness metrics for a third-party facial-analysis model and reports them in a standardised, regulator-legible format. The platform is organised into five cooperating modules covering data ingestion, model interfacing, fairness-metric computation, explainability, and report generation. Its analytical core computes four disparity metrics, demographic parity difference, equalized odds difference, equal opportunity difference, and disparate impact ratio, using AIF360 and Fairlearn as independent, cross-validating backends, with every reported disparity accompanied by a chi-squared significance test and a bootstrap confidence interval rather than a bare threshold. A SHAP-based explainability layer attributes flagged disparities to input features, distinguishing genuinely demographic effects from unrelated confounds. Findings are assembled into an exportable, Model-Cards- and Datasheets-structured report in which every metric is traced to its specific basis under Article~10 of the EU AI Act and the corresponding function of the NIST AI Risk Management Framework. The design is validated against the FairFace and UTKFace benchmark datasets, with a FairFace-trained convolutional-network baseline serving as the platform's own case study. BiasAperture is scoped strictly as diagnostic: it identifies and statistically characterises disparities but does not mitigate bias, retrain models, or generate synthetic demographic data.

\vspace{1cm}
\noindent
\textbf{Keywords:} Algorithmic Fairness, Bias Auditing, Demographic Parity, Disparate Impact, EU AI Act, Explainable AI, Facial Analysis, Fairness Metrics, NIST AI RMF, SHAP
